%================================================
% 5. TONG KET
%================================================
\section{Tổng kết}

Các Biểu đồ cho thấy tổng phát thải CO\textsubscript{2} toàn cầu có xu hướng tăng trong giai đoạn 2000--2020, với các nhịp giảm ngắn quanh 2008--2009 và 2019--2020. Khi tách theo khu vực, mức tăng không phân bố đều: East Asia \& Pacific mở rộng tỷ trọng mạnh nhất, trong khi North America và Europe \& Central Asia giảm tỷ trọng so với mốc đầu. Xét về quy mô phát thải, khoảng cách giữa các khu vực không thu hẹp mà dịch chuyển trung tâm, từ những khu vực phát triển sớm sang East Asia \& Pacific và South Asia.


Góc nhìn CO\textsubscript{2} bình quân đầu người cho một bức tranh khác với phát thải tổng. Bốn nhóm thu nhập tách tầng rõ trên Biểu đồ: High income vẫn cao nhất dù có xu hướng giảm, Upper middle income tăng rõ, còn Low income nằm sát đáy trong toàn giai đoạn. Biểu đồ khoảng cách cho thấy chênh lệch High income -- Low income thu hẹp về cuối giai đoạn nhưng vẫn lớn, và bản đồ CO\textsubscript{2}/người cho thấy các quốc gia nổi bật không trùng với nhóm phát thải tổng lớn nhất. Khi chuẩn hóa theo dân số, khoảng cách phát triển giữa các nhóm thu nhập hiện rõ hơn và vẫn duy trì.


Các Biểu đồ cường độ carbon cho thấy lượng CO\textsubscript{2} trên một đơn vị GDP giảm ở nhiều nhóm thu nhập và khu vực, và scatter plot ghi nhận một nhóm quốc gia có GDP/người tăng nhanh trong khi tổng CO\textsubscript{2} tăng chậm hơn, chỉ một số ít quốc gia giảm tổng CO\textsubscript{2} tuyệt đối. Tuy nhiên, mức cải thiện không đồng đều giữa các nhóm. Việc cường độ carbon giảm trong khi tổng phát thải toàn cầu vẫn tăng cho thấy phần cải thiện hiệu quả phát thải bị bù lại bởi quy mô kinh tế và dân số; báo cáo chỉ ghi nhận hiện tượng quan sát được trên Biểu đồ và không diễn giải thành nguyên nhân. Vì vậy, khoảng cách về hiệu quả phát thải trên một đơn vị kinh tế giữa các nhóm vẫn còn.


Các Biểu đồ năng lượng tái tạo cho thấy tỷ trọng tái tạo thay đổi không đồng nhất: Low income có tỷ trọng cao trên Biểu đồ, High income tăng từ mức thấp hơn, trong khi hai nhóm thu nhập trung bình có xu hướng giảm. Đây là tỷ trọng quan sát được trên Biểu đồ và không hàm ý hạ tầng hay chính sách năng lượng của một nhóm tốt hơn nhóm khác. Scatter plot và bubble chart cho thấy tiêu thụ năng lượng/người, CO\textsubscript{2}/người, GDP/người và tỷ trọng tái tạo không nằm trên một quan hệ tuyến tính đơn giản, nên chuyển dịch năng lượng bổ sung thêm một chiều cho bức tranh khoảng cách phát triển, bên cạnh phát thải và hiệu quả phát thải.


Tổng hợp bốn góc nhìn, Dashboard và các Biểu đồ cho thấy khoảng cách phát triển giữa các nhóm thu nhập và khu vực vẫn lớn trong giai đoạn 2000--2020 và dịch chuyển theo quy luật riêng ở từng chỉ số: trung tâm phát thải tổng dời sang East Asia \& Pacific, khoảng cách phát thải bình quân đầu người thu hẹp mà chưa biến mất, và hiệu quả phát thải trên một đơn vị kinh tế cải thiện không đồng đều. Bốn chỉ số -- tổng phát thải CO\textsubscript{2}, phát thải bình quân đầu người, cường độ carbon và tỷ trọng năng lượng tái tạo -- vì vậy bổ sung cho nhau, mỗi chỉ số làm rõ một mặt của khoảng cách phát triển. Đặt cạnh nhau trên cùng một Dashboard, các chỉ số này hợp thành một bức tranh đầy đủ và nhất quán về khoảng cách phát triển.
